\documentclass[11pt,a4paper]{article}

\usepackage[utf8]{inputenc}
\usepackage[cyr]{aeguill}
\usepackage[francais]{babel}

% Les packages
%\usepackage[french]{babel}
\usepackage{fancybox}
\usepackage{graphics}
\usepackage{color}
\usepackage{latexsym}
\usepackage{amsmath,amsthm,amssymb}
\usepackage[plain]{fullpage} 
\usepackage{verbatim}
\usepackage{times,epsfig}
\usepackage{listings}
\usepackage{multirow}
\usepackage{amsmath,amsthm,amssymb,amscd,txfonts,bbm} % RAJOUT ICI !!!
\usepackage{graphics,epsfig} % RAJOUT ICI !!!
%\usepackage{algorithmic}
\usepackage[final]{pdfpages} 
\usepackage{fancyvrb}
\usepackage{enumerate}
\usepackage{listingsutf8}
\usepackage{comment}
\usepackage{minted}

% small et verbatim: to be coherent with \file
\newenvironment{smallverbatim}%
{\endgraf\small\verbatim}{\endverbatim}

\newenvironment{qv}
{\quote\Verbatim}
{\endquote\endVerbatim}

\usepackage{lastpage}

\newcommand{\terme}[1]{\texttt{#1}}
\newcommand{\termSet}[1]{\{\texttt{#1}\}}
\newcommand{\guill}[1]{<<\,#1\,>>}
% Macro pour l'inclusion de figure  
\newcommand{\fig}[2]{%
  \begin{figure*}[hbt]
   \begin{center}
      \includegraphics[width=12cm]{#1}
  \caption{\label{#1}  #2}
  \end{center}
 \end{figure*}
}
\newcommand{\figTex}[2]{%
  \begin{figure*}[hbt]
 \centerline{\input{#1.pstex_t}}
  \caption{\label{#1}  #2}
 \end{figure*}
}

% Macro résumés
\definecolor{turquoise}{rgb}{0.20 ,0.60, 0.60}
\definecolor{vert}{rgb}{0.40,  0.60 ,0.00}
\definecolor{violet}{rgb}{0.80 , 0.00 , 0.80}
\definecolor{bleu}{rgb}{0.20,  0.20,  0.80}
\definecolor{rose}{rgb}{1.00,  0.20,  0.40}

\lstset{basicstyle=\footnotesize,showstringspaces=false,language=XML,breaklines=false,columns=fullflexible,frame=shadowbox, captionpos=b}

\newenvironment{solution}[0]{
~\\ \color{red} \textbf{Solution : }\color{black}}{%
}
%\excludecomment{solution}


\newcommand{\cle}[1]{\textbf{#1}}
\def\fk#1{\textit{#1}}
\def\tble#1{\textit{#1}}

\parindent=0pt           
                             
\usepackage{url}

\newtheorem{Exercice}{Definition}


\begin{document}

\title{Sujet d'examen UE NFE204 FOAD - Session 2}
\author{Année universitaire 2019-2020, 5 septembre}
\date{\empty}

\maketitle

{\bf \Large Consignes.}
Cet examen est exceptionnellement organisé à distance. 
Vous devez rendre votre copie numérisée à 17h00 au plus tard. 
Aucune copie ne sera acceptée après cette heure limite. 
Les seuls formats acceptés sont le PDF et le JPEG, ceci afin d'éviter tout problème de lecture.  
Prenez vos dispositions pour être en mesure de fournir le bon format le moment venu. Une activité 
de discussion en direct (\guill{\textit{chat}}) est ouverte dans l'espace Moodle pour d'éventuelles questions.


\section*{Première partie : modélisation semi-structurée (6 pts)}

Un grand magasin enregistre tous les tickets de caisse des transactions effectuées par ses clients.
Un ticket de caisse comprend la date et l'heure de chaque transaction, le numéro de la caisse, le prix total,
et la liste des articles achetés avec, pour chacun, la quantité et le prix unitaire.  Le ticket de caisse comprend également
l'identifiant du client si ce dernier a utilisé sa carte de fidélité. Sinon cet identifiant
est à \texttt{null}.


\begin{enumerate}
  \item Dans une base relationnelle, combien de tables faudrait-il pour représenter 
         les informations d'un ticket de caisse (expliquer et donner le schéma de ces tables); combien d'insertions de ligne
             faut-il faire pour chaque ticket\,?
    \item Si on choisit une base NoSQL, comment représenter un ticket de caisse en JSON.
        Donnez un exemple simple. Combien d'insertions faut-il faire\,?
   \item Je veux compter le nombre de transactions par caisse et par jour.  Quelle est la requête SQL si les données
       sont dans une base relationnelle\,?
    \item Même question si la base est NoSQL\,: comment calculer le nombre de transactions par caisse et par jour
    \item Indiquez les avantages et inconvénients de chaque solution, relationnelle ou NoSQL.
 \end{enumerate}


\begin{solution}

\begin{enumerate}
   \item En relationnel, il faut 2 tables, une pour les tickets avec les champs \texttt{noCaisse},
      \texttt{date}, \texttt{heure} et \texttt{prixTotal}, une autre pour les articles avec \texttt{idArticle},
         \texttt{quantité} et \texttt{prixUnitaire}. Un champ \texttt{idClient}
          dans la table \texttt{Ticket} peut être à \texttt{null}. La table \texttt{Article} comprend
         une clé étrangère \texttt{idTicket} pour savoir de quel ticket l'achat de l'article fait partie.
         
          Il faut une insertion dans la table \texttt{Ticket},
           et autant d'insertions dans la table \texttt{Article} qu'il y a d'articles achetés dans
           une même transaction. 
      
      \item En JSON, par exemple:

\begin{minted}[frame=leftline,
               baselinestretch=.9,
               fontsize=\footnotesize,
               framesep=2mm]{json}
{
  "_id": 978,
  "noCaisse": 12,
  "date": "01/01/2017",
  "heure": "13:30",
  "prixTotal": 58.47,
  "idClient": null,
  "articles": [
     {"idArticle": 1982, "quantité": 5,   "prix": 10.45},
     {"idArticle": 34, "quantité": 1,   "prix": 45},
     {"idArticle": 154, "quantité": 3,   "prix": 3.02},
  ]
}
\end{minted}

Une insertion suffit.

\item En SQL:

\begin{minted}{sql}
select count(*) from Ticket t, Article a 
where a.idTicket = t.id
group by date, noCaisse
\end{minted}

\item En NoSQL, il faut écrire et exécuter un programme MapReduce.

La fonction de Map émet  une paire intermédiaire dont la clé est le critère de regroupement,
soit la paire (date, numéro de caisse).

\begin{minted}{bash}
function fonctionMap ($ticket)
{
  emit ( [$ticket.date, $ticket.noCaisse], 1)
}
\end{minted}

La fonction de Reduce effectue le compte des paires intermédiaires, autrement dit des tickets
ayant même  date et même numéro de caisse.

\begin{minted}{bash}
function fonctionReduce ($clé, $valeurs)
{
  return ($clé, count($valeurs))
}
\end{minted}
\end{enumerate}
\end{solution}


\section*{Deuxième partie : MapReduce (4 pts)}

Je veux calculer le montant moyen des tickets par client, pour les clients dont
je connais l'identifiant. 

\begin{enumerate}
 \item Décrivez la modélisation MapReduce de ce calcul. 
 Donnez le pseudo-code de chaque
fonction, ou indiquez par un texte clair son déroulement.
\item Donnez la forme de la chaîne de traitement (\textit{workflow}), avec des opérateurs Pig ou Spark,
      qui implante ce calcul.
\end{enumerate}


\begin{solution}

\begin{enumerate}
 
\item La fonction de Map  ne prend en compte que les tickets ayant un numéro de client non nul,
   et émet  une paire intermédiaire dont la clé est le critère de regroupement,
soit le numéro de client, et la valeur le prix total du ticket.

\begin{minted}{bash}
function fonctionMap ($ticket)
{
  if ($ticket.idClient not null) then
      emit ( $ticket.idClient, $ticket.prixTotal)
  end if
}
\end{minted}

La fonction de Reduce effectue la moyenne des prix.
\begin{minted}{bash}
function fonctionReduce ($idClient, $listePrix)
{
  return ($idClient, avg($listePrix))
}
\end{minted}

\item En Pig:

\begin{minted}{pig}
tickets_client = filter tickets BY idClient not null;
tickets_groupes = group tickets_client by idClient;
avg_par_client = foreach tickets_groupes generate group, AVG(tickets_groupes.prixTotal);
\end{minted}
\end{enumerate}
\end{solution}

\section*{Troisième partie : recherche d'information (4 pts)}

Les articles sont décrits par des mots-clé.
Voici les mots-clé de 4 articles $a_1$, $a_2$, $a_3$, $a_4$.

\begin{enumerate}
  \item laitage, céréale, fruit, légume.
  \item céréale, viande, protéine, sucre, gras 
  \item laitage, fruit, sucre, amidon, vitamine
  \item céréale, viande

\end{enumerate}

\smallskip

\begin{enumerate}
   \item Quel est l'idf de "céréale"\,? Quelle interprétation donnez-vous à ce résultat\,?

\begin{solution}
L'idf est proche de 0 ($\log (4/3) = 0,12$), ce qui signifie que le terme est très 
courant et peu discriminant pour une recherche.
\end{solution}

  \item Une recherche avec le mot "viande" ramènera quels documents, et dans quel ordre (pas la peine de faire des calculs).
    
\begin{solution}
Le résultat est, dans l'ordre,  $a_4$ puis $a_2$. Les deux parlent de viande, mais l'importance relative
du terme est plus grande pour $a_4$. Pour le calcul, la norme de  $a_4$ est plus petite, et la division par
la norme donnera un résultat plus grand.
\end{solution}

  \item Quelle est la similarité cosinus entre les articles $a_1$ et $a_3$\,?
\begin{solution}
La norme de $a_1$ est $\sqrt{4} = 2$. La norme de $a_3$ est $\sqrt{5}$. Le cosinus est

$$\frac{1 + 1}{2 \times \sqrt{5}} = \frac{1}{\sqrt{5}}$$ 
\end{solution}
  
  \end{enumerate}

\section*{Quatrième partie : systèmes distribués (2 pts)}

Un groupe de serveurs en mode maître-esclave comprend 5 n{\oe}uds. Suite à une fragmentation
du réseau, le maître se retrouve  avec un seul esclave.  Expliquez ce qui se passe alors.

\section*{Cinquième partie : questions de cours (4 pts)}


\begin{enumerate}
   \item Donnez une définition de la scalabilité.
   \item Définissez le partitionnement d'une collection dans un système distribué. Quel est l'intérêt?
   \item Pour un traitement, est-il en général plus efficace de lire les données sur le disque local
      ou dans la mémoire RAM d'un serveur de la même armoire\,?
    \item Quelle sont les opérations proposées par une ressource dans une architeture RESTful\,?
\end{enumerate}


\end{document}
